\documentclass[11pt,a4paper]{article}

\usepackage[margin=2.35cm,headheight=15pt]{geometry}
\usepackage{fontspec}
\usepackage{xeCJK}
\setCJKmainfont{Noto Serif CJK SC}[AutoFakeSlant=true]
\setCJKsansfont{Noto Sans CJK SC}[AutoFakeSlant=true]
\setCJKmonofont{Noto Sans CJK SC}
\xeCJKsetup{PunctStyle=kaiming}

\usepackage{amsmath,amssymb,amsthm,mathtools,bm}
\usepackage{booktabs}
\usepackage{enumitem}
\usepackage{xcolor}
\usepackage{tikz}
\usetikzlibrary{arrows.meta,positioning,calc,decorations.pathreplacing,decorations.markings}
\usepackage[most]{tcolorbox}
\usepackage{hyperref}
\usepackage{fancyhdr}
\usepackage{titlesec}
\usepackage{array}
\usepackage{microtype}

\hypersetup{
  colorlinks=true,
  linkcolor=blue!50!black,
  urlcolor=blue!60!black,
  citecolor=blue!50!black
}

\definecolor{syblue}{RGB}{70,30,110}
\definecolor{sysoft}{RGB}{242,236,250}
\definecolor{warnbg}{RGB}{255,246,230}
\definecolor{okbg}{RGB}{232,246,236}

\theoremstyle{definition}
\newtheorem{definition}{定义}[section]
\newtheorem{example}{例子}[section]
\theoremstyle{plain}
\newtheorem{theorem}{定理}[section]
\newtheorem{proposition}{命题}[section]
\theoremstyle{remark}
\newtheorem{remark}{注记}[section]

\newtcolorbox{keybox}[1][]{
  colback=sysoft,colframe=syblue,fonttitle=\bfseries,
  title={核心要点},#1
}
\newtcolorbox{warnbox}[1][]{
  colback=warnbg,colframe=orange!70!black,fonttitle=\bfseries,
  title={容易混淆},#1
}
\newtcolorbox{hearbox}[1][]{
  colback=okbg,colframe=green!40!black,fonttitle=\bfseries,
  title={听课提示},#1
}

\pagestyle{fancy}
\fancyhf{}
\lhead{复旦粒子物理暑期学校 2026}
\rhead{Generalized Symmetry / Daniel Brennan}
\cfoot{\thepage}
\titleformat{\section}{\Large\bfseries\color{syblue}}{\thesection}{0.8em}{}
\titleformat{\subsection}{\large\bfseries\color{syblue!80!black}}{\thesubsection}{0.7em}{}

\newcommand{\dd}{\mathrm{d}}
\newcommand{\Dd}{\mathrm{D}}
\DeclareMathOperator{\tr}{tr}

\begin{document}

\begin{center}
{\LARGE\bfseries 广义对称性课前讲义}\\[0.45em]
{\large 复旦粒子物理暑期学校 2026}\\[0.25em]
{\large Generalized Symmetry \quad T.\ Daniel Brennan (University of Birmingham)}\\[0.8em]
{\normalsize 课前预习稿　从普通对称性写起　推导写全　英文术语单列}\\[0.35em]
{\small 依据 Indico 课表：8 月 17--19 日共四讲。本讲义非正式课堂讲义。}
\end{center}

\tableofcontents
\newpage

%======================================================================
\section{如何使用本讲义}
%======================================================================

广义对称性（generalized global symmetry）把``对称性 = 群 $G$ 作用在粒子上''这句话拆开，改写成：

\begin{quote}
\textbf{对称性由拓扑算符实现。}算符的空间维数决定它作用在什么对象上。
\end{quote}

Daniel Brennan 本人有一份专为高能现象学读者写的综述：
\emph{Introduction to Generalized Global Symmetries in QFT and Particle Physics}, arXiv:2306.00912。
课堂大概率沿这条路线：连续高形式对称性 $\to$ 离散对称与中心对称 $\to$ 高群 $\to$ 非可逆对称性。本讲义把综述默认已经会的``普通 QFT 对称性''全部补上，并把 Maxwell 理论当作贯穿始终的例子。

四讲课表（Asia/Shanghai）：

\begin{center}
\begin{tabular}{lll}
\toprule
日期 & 时间 & 内容 \\
\midrule
8 月 17 日（一） & 13:30--14:30 & Generalized Symmetry \\
8 月 18 日（二） & 13:30--14:30 & Generalized Symmetry \\
8 月 19 日（三） & 09:20--10:20, 10:40--11:40 & Generalized Symmetry \\
\bottomrule
\end{tabular}
\end{center}

\begin{hearbox}
Brennan 的综述明确\textbf{不用范畴论}。预习时不要先去抱着 fusion category 读。先会：Noether 流、Wilson 线、拓扑算符与 linking、Maxwell 的电/磁 1-形式对称性。范畴语言若出现，当作可选包装。
\end{hearbox}

%======================================================================
\section{预备知识}
\label{sec:prereq}
%======================================================================

\subsection{经典场论}

时空 $d$ 维（粒子物理默认 $d=4$），场 $\phi(x)$，作用量
\begin{equation}
  S[\phi]=\int \dd^d x\,\mathcal{L}(\phi,\partial\phi).
\end{equation}
运动方程由 $\delta S=0$ 得到 Euler--Lagrange 方程
\begin{equation}
\label{eq:EL}
  \frac{\partial\mathcal{L}}{\partial\phi}-\partial_\mu\frac{\partial\mathcal{L}}{\partial(\partial_\mu\phi)}=0.
\end{equation}

\subsection{微分形式（最小量）}

$p$-形式是完全反对称的 $(0,p)$ 张量。外导数 $\dd$ 满足 $\dd^2=0$。Hodge 星 $*$ 把 $p$-形式映到 $(d-p)$-形式。对 1-形式流 $j=j_\mu\dd x^\mu$，
\begin{equation}
  \partial_\mu j^\mu=0 \quad\Longleftrightarrow\quad \dd{*}j=0.
\end{equation}
Maxwell 场强 $F=\dd A$ 是 2-形式；Bianchi 恒等式 $\dd F=0$ 就是 $\partial_{[\lambda}F_{\mu\nu]}=0$。

积分：把 $p$-形式在 $p$ 维子流形上积分。Stokes 公式
\begin{equation}
  \int_\Sigma \dd\omega=\int_{\partial\Sigma}\omega.
\end{equation}
后面``把流在余维度 $1$ 的流形上积分得到荷''，靠的就是这句话。

\subsection{群与表示}

普通全局对称群 $G$ 作用在局域算符上：
\begin{equation}
  \mathcal{O}(x)\ \mapsto\ g\cdot\mathcal{O}(x),\qquad g\in G.
\end{equation}
$U(1)$ 的表示由整数电荷 $q$ 标记：$e^{i\alpha}\cdot\mathcal{O}=e^{iq\alpha}\mathcal{O}$。非阿贝尔群有换位子，高形式对称性（$p\ge 1$）却必须是阿贝尔的---后面会看到为什么。

\subsection{QFT 中的算符}

\begin{itemize}[leftmargin=1.4em]
  \item \textbf{局域算符} $\mathcal{O}(x)$：插入在点上，激发粒子。
  \item \textbf{线算符}（line operator）：插入在一维曲线上。Wilson 线是最重要的例子。
  \item \textbf{表面算符}、更高维的缺陷（defect）。
\end{itemize}
关联函数
\begin{equation}
  \langle \mathcal{O}_1(x_1)\cdots\mathcal{O}_n(x_n)\rangle
  =\frac{1}{Z}\int \mathcal{D}\phi\, e^{-S[\phi]}\,\mathcal{O}_1\cdots\mathcal{O}_n
\end{equation}
（欧氏记号）。对称性最干净的定义，是这些算符如何变换，以及哪些算符的插入不依赖于连续形变。

%======================================================================
\section{普通对称性与 Noether 定理}
\label{sec:noether}
%======================================================================

\subsection{有限变换与无穷小变换}

设连续对称群 $G$ 的无穷小参数为常值 $\alpha$，场变分
\begin{equation}
  \phi(x)\ \mapsto\ \phi(x)+\alpha\,\delta\phi(x).
\end{equation}
``是对称性''意味着：当 $\alpha$ 为常数时 $\delta S=0$（或只差一个边界项）。

\subsection{完整推导：位置相关的参数}

把 $\alpha$ 提升为 $\alpha(x)$。此时作用量一般不再不变，但变分只能靠 $\partial\alpha$，因为 $\alpha$ 为常数时变分为零。于是必然有
\begin{equation}
\label{eq:noether-var}
  \delta S=\int \dd^d x\,(\partial_\mu\alpha)\, j^\mu
  =-\int \dd^d x\,\alpha\,\partial_\mu j^\mu
  \quad\text{（丢弃边界项）}.
\end{equation}
在运动方程上 $\delta S=0$ 对任意 $\alpha(x)$ 成立，故
\begin{equation}
\label{eq:cons}
  \partial_\mu j^\mu=0.
\end{equation}
这就是 Noether 定理：连续整体对称性 $\Rightarrow$ 守恒流。

\begin{example}[复标量的 $U(1)$]
$\mathcal{L}=\partial_\mu\phi^\dagger\partial^\mu\phi-V(\phi^\dagger\phi)$，$\phi\mapsto e^{i\alpha}\phi$。无穷小 $\delta\phi=i\phi$，$\delta\phi^\dagger=-i\phi^\dagger$。标准计算给出
\begin{equation}
  j^\mu=i\bigl(\phi^\dagger\partial^\mu\phi-\phi\,\partial^\mu\phi^\dagger\bigr),\qquad
  Q=\int \dd^{d-1}x\, j^0.
\end{equation}
$Q$ 计的是粒子数减去反粒子数。
\end{example}

\begin{example}[时空平移]
$\delta\phi=-\alpha^\nu\partial_\nu\phi$ 给出能量动量张量 $T^{\mu\nu}$，$\partial_\mu T^{\mu\nu}=0$。能量是 $0$-形式对称性对应的荷，对象仍是局域算符（粒子）。
\end{example}

\subsection{荷作为算符}

量子理论中，
\begin{equation}
  Q=\int_{\Sigma_{d-1}} {*j},
\end{equation}
$\Sigma_{d-1}$ 是一个类空切片。守恒性 $\dd{*}j=0$ 加上 Stokes，使得只要把 $\Sigma$ 连续形变且不碰到带电算符，$Q$ 不变。于是
\begin{equation}
  U(g)=e^{i\alpha Q}
\end{equation}
满足
\begin{equation}
\label{eq:action0}
  U(g)\,\mathcal{O}(x)\,U(g)^{-1}=g\cdot\mathcal{O}(x).
\end{equation}
若 $\mathcal{O}$ 电荷为 $q$，则 $U(e^{i\alpha})\mathcal{O}U^{-1}=e^{iq\alpha}\mathcal{O}$。

\subsection{Ward 恒等式}

把~\eqref{eq:noether-var} 放到路径积分里，得到插入守恒流的关联函数恒等式。对 $U(1)$，
\begin{equation}
  \partial_\mu\langle j^\mu(x)\,\mathcal{O}_1(x_1)\cdots\rangle
  =\sum_i q_i\,\delta^{(d)}(x-x_i)\,\langle \mathcal{O}_1\cdots\rangle.
\end{equation}
这是``对称性作用在局域算符上''的关联函数版。

%======================================================================
\section{规范对称还是全局对称}
\label{sec:gauge-vs-global}
%======================================================================

\begin{definition}[全局 vs.\ 规范]
全局对称的参数 $\alpha$ 是时空上的常数（或至少在无穷远处不能任意规范掉）。规范``对称''的参数 $\alpha(x)$ 是局域的冗余：两个场位形若只差一个规范变换，代表同一物理态。
\end{definition}

\begin{keybox}
现代观点：\textbf{规范对称不是物理对称，而是描述的冗余}。真正能把态分类、能给出选择定则、能自发破缺的，是\textbf{全局对称}。广义对称性讨论的全部是全局对称；尽管它们常常用规范场的拓扑对象（Wilson 线、't~Hooft 线）来探测。
\end{keybox}

把全局 $U(1)$ 耦合到\textbf{背景}规范场 $A_\mu$（不积分 $A$，只把它当作探针）：
\begin{equation}
  S\ \supset\ \int A_\mu j^\mu.
\end{equation}
背景场规范变换 $A\mapsto A+\dd\lambda$ 下，作用量不变当且仅当 $\dd{*}j=0$。若不变性失败，失败项就是 't~Hooft 反常。Brennan 会反复使用``打开背景规范场''这一手法。

%======================================================================
\section{拓扑算符：对称性的现代定义}
\label{sec:topo}
%======================================================================

\subsection{从切片上的荷到任意封闭流形}

普通荷 $Q$ 定义在类空切片上。但 $\dd{*}j=0$ 允许把积分流形换成任意封闭的 $(d-1)$ 维流形 $M_{d-1}$：
\begin{equation}
\label{eq:sdo}
  U_g(M_{d-1})=\exp\Bigl(i\alpha\int_{M_{d-1}}{*j}\Bigr),\qquad g=e^{i\alpha}\in U(1).
\end{equation}
$U_g(M_{d-1})$ 称为\textbf{对称性缺陷算符}（symmetry defect operator, SDO）。

\subsection{为什么叫``拓扑''}

若 $M$ 连续形变成 $M'$，且形变扫过的区域里没有带电算符，则 Stokes 加流守恒给出
\begin{equation}
  U_g(M)=U_g(M').
\end{equation}
关联函数只依赖于 $M$ 的拓扑类（以及它与带电算符的 linking），不依赖于具体几何形状。这样的算符叫做拓扑算符。

\begin{figure}[h]
\centering
\begin{tikzpicture}[scale=1.05]
  \draw[thick] (0,0) circle (1.3);
  \fill (0,0) circle (2pt) node[below right] {$\mathcal{O}_q$};
  \node at (0,1.55) {$U_g(S^{d-1})$};
  \draw[-{Stealth},thick] (3.2,0) -- (4.3,0);
  \draw[thick] (7.1,0.15) ellipse (1.5 and 1.05);
  \fill (7.1,0.15) circle (2pt) node[below right] {$\mathcal{O}_q$};
  \node at (7.1,1.45) {$U_g(M')$};
  \node[font=\small] at (2.15,-1.7) {连续形变不改变关联函数};
\end{tikzpicture}
\caption{$U_g$ 包围带电局域算符。把它捏成任意形状，只要仍链绕一次，相位 $g^q$ 不变。}
\end{figure}

\subsection{作用 = linking}

把 $U_g(S^{d-1})$ 从无穷远处收缩，直到它包住 $\mathcal{O}_q(x)$。由~\eqref{eq:action0}，
\begin{equation}
\label{eq:link0}
  U_g(S^{d-1})\ \mathcal{O}_q(x)\ =\ g^{q}\ \mathcal{O}_q(x).
\end{equation}
更拓扑的说法：对称性缺陷与带电算符的 linking 数决定相位。

\begin{definition}[0-形式全局对称性]
$d$ 维 QFT 的 0-形式对称群 $G$ 由一族拓扑算符 $\{U_g(M_{d-1})\}_{g\in G}$ 给出，它们：
\begin{enumerate}[leftmargin=1.6em]
  \item 支持在余维度 $1$ 的封闭流形上；
  \item 融合规则为 $U_{g_1}U_{g_2}=U_{g_1g_2}$；
  \item 作用在局域（0 维）算符上，如~\eqref{eq:link0}。
\end{enumerate}
\end{definition}

连续对称只是特例：拓扑算符可由守恒流积分构造。离散对称没有 Noether 流，但仍然可以有拓扑算符。这一点是走向``广义''的关键---\textbf{流不是本质，拓扑算符才是}。

%======================================================================
\section{从 0-形式到 $p$-形式}
\label{sec:pform}
%======================================================================

普通对称性作用在粒子（局域算符）上。若荷的载体不是粒子而是弦、膜呢？

\begin{definition}[$p$-形式全局对称性]
$d$ 维 QFT 具有群 $G$ 的 $p$-形式全局对称，若存在拓扑算符 $U_g(M_{d-p-1})$，支持在余维度 $p+1$ 的封闭流形上，融合服从群乘法，并作用在 $p$ 维带电算符 $W(C_p)$ 上：
\begin{equation}
\label{eq:linkp}
  U_g(M_{d-p-1})\,W(C_p)= g^{q[C]}\,W(C_p),
\end{equation}
其中 $q[C]$ 由 linking 数给出。
\end{definition}

对照表：

\begin{center}
\begin{tabular}{cccc}
\toprule
$p$ & 带电对象 & 对称性缺陷的余维 & 连续情形的流 \\
\midrule
$0$ & 局域算符 / 粒子 & $1$（面） & 1-形式流 $j_1$，$\dd{*}j_1=0$ \\
$1$ & 线算符 / 弦 & $2$ & 2-形式流 $j_2$，$\dd{*}j_2=0$ \\
$p$ & $p$ 维算符 & $p+1$ & $(p+1)$-形式流，$\dd{*}j_{p+1}=0$ \\
\bottomrule
\end{tabular}
\end{center}

连续 $p$-形式对称的 Noether 流是 $(p+1)$-形式 $j_{p+1}$，守恒
\begin{equation}
  \dd{*}j_{p+1}=0,
\end{equation}
背景场是 $(p+1)$-形式 $A_{p+1}$，耦合 $\int A_{p+1}\wedge{*}j_{p+1}$。背景规范变换
\begin{equation}
  A_{p+1}\mapsto A_{p+1}+\dd\lambda_p.
\end{equation}

\subsection{为什么 $p\ge 1$ 必须阿贝尔}

两个余维度 $\ge 2$ 的缺陷可以互相绕过对方而不相交。因此 $U_{g}U_{h}$ 与 $U_{h}U_{g}$ 可以通过形变互换，故 $G$ 必须交换。$0$-形式缺陷是余维度 $1$ 的``墙''，不能随便穿过，所以 $0$-形式群可以非阿贝尔。

\begin{warnbox}
``$p$-形式''中的 $p$ 是\textbf{带电算符的维数}，不是流的形式次数。流是 $(p+1)$-形式。课堂若写 $G^{(p)}$，括号里的 $p$ 也是带电对象的维数。
\end{warnbox}

%======================================================================
\section{典型例子：Maxwell 理论}
\label{sec:maxwell}
%======================================================================

这是整门课最重要的例子，必须能独立复述。

\subsection{作用量与方程}

无物质的 $U(1)$ 规范理论（欧氏或闵氏，这里用形式语言）
\begin{equation}
  S=-\frac{1}{2e^2}\int F\wedge{*}F,\qquad F=\dd A.
\end{equation}
运动方程
\begin{equation}
\label{eq:eom}
  \dd{*}F=0,
\end{equation}
Bianchi 恒等式
\begin{equation}
\label{eq:bianchi}
  \dd F=0.
\end{equation}

\subsection{电 1-形式对称性}

\eqref{eq:eom} 正是 $\dd{*}j_e=0$，其中电 2-形式流
\begin{equation}
  j_e=\frac{1}{e^2}F\qquad\text{（约定可差常数）}.
\end{equation}
因此存在连续 $U(1)^{(1)}_e$ 电 1-形式对称。对称性缺陷是在封闭 2-面 $\Sigma_2$ 上积分 $*F$：
\begin{equation}
  U_e^{(\alpha)}(\Sigma_2)=\exp\Bigl(\frac{i\alpha}{e^2}\int_{\Sigma_2}{*F}\Bigr).
\end{equation}
它作用的带电对象是 Wilson 线
\begin{equation}
  W_q(C)=\exp\Bigl(iq\int_C A\Bigr).
\end{equation}
物理图像：Wilson 线是一条无穷重的探针电荷的世界线。把一张面 $\Sigma_2$ 链绕 $C$ 一次，Gauss 定律给出相位 $e^{iq\alpha}$。

\subsection{磁 1-形式对称性}

Bianchi 恒等式~\eqref{eq:bianchi} 本身也是守恒律：磁流 $j_m={*}F$ 的对偶说法是 $\dd F=0$。缺陷
\begin{equation}
  U_m^{(\beta)}(\Sigma_2)=\exp\Bigl(i\beta\int_{\Sigma_2} F\Bigr)
\end{equation}
作用在 't~Hooft 线 $T(C')$ 上。't~Hooft 线对应磁单极探针。电、磁 1-形式对称在自由 Maxwell 中同时存在。

\subsection{选择定则}

$U(1)^{(1)}_e$ 禁止 Wilson 线``断掉''：电荷线必须延伸到无穷远或形成闭合圈。这就是普通电荷守恒的 1-形式版---守恒的不是粒子数，而是通量线。

\begin{example}[紧致标量 / 2d 光子]
$d=2$ 的紧致标量等价于 2d 光子。$0$-形式与 $1$-形式对称互换角色。课堂可能用这个低维例子先建立直觉，再回到 4d Maxwell。
\end{example}

%======================================================================
\section{Higgsing、禁闭与 Landau 范式的扩充}
\label{sec:phases}
%======================================================================

Landau 范式：用全局对称是否自发破缺来标记相。高形式对称把这句话扩展到规范理论的相。

\subsection{Higgs 相：电 1-形式破缺}

加入电荷为 $1$ 的物质场后，Wilson 线可以在物质上终止，电 1-形式对称被显式破坏（或被 Higgs 机制破缺）。Wilson 线服从周长律（perimeter law）：其期望值随长度指数衰减，但不是面积律。直观上，通量管被电荷屏蔽。

\subsection{禁闭相：电 1-形式未破缺}

纯 Yang--Mills 的禁闭相中，Wilson 圈服从面积律
\begin{equation}
  \langle W(C)\rangle\sim e^{-\sigma\cdot\mathrm{Area}(C)}.
\end{equation}
这表示有一张物理的色通量弦张在回路上。电 1-形式（更精确地：中心 1-形式）\textbf{未破缺}：弦不能断。't~Hooft 线则被破缺。

\subsection{库仑相}

自由光子：Wilson 与 't~Hooft 都是周长律，电、磁 1-形式都未破缺（连续对称；注意 Coleman--Mermin--Wagner 对高形式对称的修正，见下）。

\begin{keybox}
\textbf{禁闭 = 电 1-形式对称未破缺；Higgs = 电 1-形式破缺。}这是广义对称性对现象学最直接的礼物：禁闭不再只是``说不清楚的非微扰现象''，而有一个对称性判据。
\end{keybox}

\subsection{高形式的 Coleman--Mermin--Wagner}

连续 $p$-形式对称在时空维数 $d\le p+2$ 时不能自发破缺。普通 $p=0$ 的情形就是 $d\le 2$ 不能破缺连续全局对称。所以 3d 光子的连续电 1-形式不会以普通 Goldstone 方式破缺。听课时若出现这句话，把它对照普通 CMW 即可。

%======================================================================
\section{离散对称、中心对称与 BF 理论}
\label{sec:discrete}
%======================================================================

Brennan 综述的第 3 节会相当技术化。预习只需抓住物理结论。

\subsection{没有 Noether 流也能有拓扑算符}

$Z_N$ 对称没有连续参数，不能靠 $\partial_\mu j^\mu=0$ 定义。但仍可存在满足
\begin{equation}
  U_g^N=1,\qquad U_g U_h=U_{gh}
\end{equation}
的拓扑算符。构造方法之一：先写一个连续 $U(1)$，再用电荷为 $N$ 的物质把它破到 $Z_N$；或直接用 BF 理论。

\subsection{BF 理论（直观）}

4d 中 BF 作用量粗略为
\begin{equation}
  S=\frac{iN}{2\pi}\int B_2\wedge \dd A_1.
\end{equation}
它是 $Z_N$ 规范理论的一种拓扑场论描述，并携带离散 0-形式与 1-形式全局对称。不必会算全部谱；要记住：\textbf{离散高形式对称的局域有效描述常常是 BF 型拓扑项}。

\subsection{Yang--Mills 的中心对称}

$SU(N)$ 纯规范理论有 $Z_N$ 中心 1-形式对称。原因：Wilson 线的表示可以相差中心元 $e^{2\pi i k/N}\mathbf{1}$，而胶子（伴随表示）看不见中心。于是有一套 $Z_N$ 的拓扑面算符，作用在基本表示 Wilson 线上。

\begin{itemize}[leftmargin=1.4em]
  \item 禁闭相：中心 1-形式未破缺，基本 Wilson 圈面积律。
  \item Higgs / 退禁闭：中心对称破缺。
\end{itemize}
若加入基本表示夸克，Wilson 线可被夸克断开，中心对称被破坏。这就是``为什么 QCD 有夸克时中心对称不是严格对称''。

\begin{example}[$SU(2)$ vs.\ $SO(3)$]
规范群是 $SU(2)$ 还是 $SO(3)=SU(2)/Z_2$，局域动力学看起来相似，但线算符谱与 1-形式对称不同。课堂很可能用这个例子强调：\textbf{广义对称性能区分局域拉氏量看不出的全局结构}。
\end{example}

%======================================================================
\section{'t Hooft 反常（预习级）}
\label{sec:anomaly}
%======================================================================

't~Hooft 反常：全局对称即使不规范，也不能被任何短距动力学破坏；它必须在 RG 流的两端匹配。操作定义：打开背景规范场后，配分函数不是背景规范不变的，
\begin{equation}
  Z[A+\dd\lambda]=e^{i\mathcal{A}[A,\lambda]}Z[A].
\end{equation}
相位 $\mathcal{A}$ 就是反常。

高形式对称同样可以有反常，包括：
\begin{itemize}[leftmargin=1.4em]
  \item 同一 $p$-形式对称自己的反常；
  \item $0$-形式与 $1$-形式之间的混合反常。
\end{itemize}
混合反常会约束 IR：例如某些相不能同时保持两个对称。Yang--Mills 中时间反演与中心对称的混合反常，是 Brennan 综述里的旗舰例子之一。

\begin{hearbox}
听到 anomaly matching 时，把它想成``IR 不能偷偷丢掉一个带反常的对称，除非它同时出现相应的 IR 自由度（Goldstone、TQFT、或其他对称破缺）''。不必在课前学会下降（inflow）的全部微分几何。
\end{hearbox}

%======================================================================
\section{高群与非可逆对称性（听课预备）}
\label{sec:beyond}
%======================================================================

\subsection{2-群}

有时 $0$-形式对称与 $1$-形式对称不能分开写，背景场变换混在一起：
\begin{equation}
  A_1\mapsto A_1+\dd\lambda_0+\cdots,\qquad
  B_2\mapsto B_2+\dd\lambda_1+\frac{\kappa}{2\pi}\lambda_0\,F_{A}+\cdots.
\end{equation}
这种结构叫 2-群（2-group）。QED 里光子与带电费米子的某些整体结构会出现它。物理后果：RG 流上，相关对称必须``整组''出现或整组消失，不能只保留其中一半。

\subsection{非可逆对称性}

普通群元有逆：$g^{-1}g=1$。非可逆对称的拓扑算符融合不再是群乘法，而可能是
\begin{equation}
  D\times D=\sum_i D_i
\end{equation}
这种``多通道''规则。4d 中，Abel 手征反常可以把连续 $U(1)$ 手征对称变成非可逆对称：经典上有 $U(1)$，量子上瞬时子破坏了它，但仍留下一套非可逆缺陷。轴子电动力学里也有类似结构。

\begin{keybox}
预习到这里即可。若课堂冲到非可逆对称，抓住一句：\textbf{拓扑算符仍然在，只是融合不再可逆}。不要卡在范畴记号上。
\end{keybox}

%======================================================================
\section{课堂可能出现的公式与术语}
\label{sec:classroom}
%======================================================================

\begin{align}
  \delta S&=\int (\partial_\mu\alpha)\,j^\mu
  && \text{Noether}\\
  \dd{*}j_{p+1}&=0
  && \text{$p$-形式守恒}\\
  U_g(M_{d-p-1})
  &&& \text{对称性缺陷}\\
  U_g\,W&=g^q\,W
  && \text{linking / 电荷}\\
  \dd{*}F&=0,\quad \dd F=0
  && \text{Maxwell 电/磁守恒}\\
  W(C)&=\exp\bigl(iq\textstyle\int_C A\bigr)
  && \text{Wilson 线}\\
  \langle W(C)\rangle&\sim e^{-\sigma\mathrm{Area}}
  && \text{面积律 = 禁闭}
\end{align}

高频词：
\textit{Noether current, topological operator, symmetry defect, linking, $p$-form symmetry, Wilson line, 't Hooft line, center symmetry, BF theory, Higgsing, confinement, 't Hooft anomaly, 2-group, non-invertible symmetry.}

%======================================================================
\section{常见困惑}
\label{sec:faq}
%======================================================================

\begin{enumerate}[leftmargin=1.8em]
  \item \textbf{把规范对称当成全局对称。}\\
  规范变换是冗余。Wilson 线的电荷来自\textbf{全局} 1-形式对称，不是来自``规范群是 $U(1)$''这句话本身。

  \item \textbf{$p$-形式的 $p$ 搞反。}\\
  $p$ = 带电算符维数。缺陷余维 = $p+1$。流是 $(p+1)$-形式。

  \item \textbf{``没有 Noether 流就没有对称性。''}\\
  离散对称没有流，仍有拓扑算符。广义对称性的定义以算符为准。

  \item \textbf{Wilson 线就是传播子。}\\
  Wilson 线是插入在路径上的算符（探针电荷历史），不是粒子传播子。它的真空期望值诊断相。

  \item \textbf{禁闭没有对称性语言。}\\
  纯规范理论中，禁闭对应中心 1-形式未破缺。有基本物质时该对称被破坏，禁闭的判据要改口。

  \item \textbf{高形式对称可以非阿贝尔。}\\
  $p\ge 1$ 时拓扑缺陷可交换，群必须阿贝尔。非阿贝尔结构出现在 $0$-形式，或在更高的范畴/高群数据里。

  \item \textbf{反常 = 对称性不存在。}\\
  't~Hooft 反常是``不能被规范、但作为全局对称仍然在，且 IR 必须匹配''。规范反常才是必须消掉的一致性条件。

  \item \textbf{非可逆 = 没有对称性。}\\
  仍然给出选择定则与 RG 约束，只是缺陷没有群意义下的逆。
\end{enumerate}

%======================================================================
\section{听课重点}
\label{sec:listen}
%======================================================================

\begin{enumerate}[leftmargin=1.8em]
  \item \textbf{第一讲（17 日下午）}：老师如何从普通 $0$-形式过渡到拓扑算符。对照本讲义第~\ref{sec:topo}--\ref{sec:pform}~节。Maxwell 例子出现时，自己默写 $\dd{*}F=0$ 与 $\dd F=0$。
  \item \textbf{第二讲（18 日下午）}：1-形式对称的 SSB、Higgs vs.\ 禁闭、可能进入离散对称。
  \item \textbf{第三、四讲（19 日上午）}：中心对称、$SU(N)$ vs.\ $SU(N)/Z_k$、反常、高群或非可逆。现象学例子（轴子、QED、Yang--Mills 相）优先记结论。
  \item 每出现一个对称，填表：\\
  \emph{$p=$？群是连续还是离散？带电算符是什么？破缺意味着哪个相？有没有反常？}
  \item 讨论时段优先问与粒子物理直接相关的：QCD 有夸克后还剩什么 1-形式结构？轴子模型里的非可逆对称有何唯象后果？
\end{enumerate}

\begin{hearbox}
若老师在黑板上画一个圈（Wilson）和一张链绕它的面（缺陷），那就是整门课的几何内核。把 linking 相位 $e^{iq\alpha}$ 写在图上。后面所有推广只是把圈换成更高维的膜、把 $U(1)$ 换成 $Z_N$ 或非可逆缺陷。
\end{hearbox}

%======================================================================
\section{进一步学习路线}
\label{sec:next}
%======================================================================

\subsection{课前}

\begin{enumerate}[leftmargin=1.6em]
  \item 从任何一本 QFT 教材复习 Noether 定理与 Ward 恒等式（Peskin--Schroeder 第 2、9 章量级）。
  \item 默写 Maxwell 方程的微分形式版，解释为何有两个 $U(1)^{(1)}$。
  \item 浏览 Brennan \& Hong, arXiv:2306.00912 的第 1--2 节（不必一次读完离散部分）。
\end{enumerate}

\subsection{课堂同步}

Gaiotto, Kapustin, Seiberg, Willett, \emph{Generalized Global Symmetries}, JHEP 2015（开山之作，可读引言与 Maxwell 节）。

\subsection{课后}

\begin{itemize}[leftmargin=1.4em]
  \item 完整读完 arXiv:2306.00912 第 3--5 节：离散、高群、非可逆。
  \item 补充：C\'ordova, Dumitrescu, Intriligator 等 TASI / 综述；Jena lectures arXiv:2407.20815（更偏应用、少范畴）。
  \item 若对 QCD 相结构感兴趣：中心对称与 Polyakov 圈的格点故事。
\end{itemize}

%======================================================================
\section{英文术语表}
\label{sec:glossary}
%======================================================================

\begin{center}
\begin{tabular}{>{\raggedright\arraybackslash}p{0.48\textwidth} p{0.44\textwidth}}
\toprule
\textbf{English} & \textbf{中文} \\
\midrule
global / gauge symmetry & 全局 / 规范对称 \\
Noether's theorem; conserved current & 诺特定理；守恒流 \\
Ward identity & 沃德恒等式 \\
local / line / surface operator & 局域 / 线 / 面算符 \\
topological operator & 拓扑算符 \\
symmetry defect operator (SDO) & 对称性缺陷算符 \\
linking number & 链绕数 \\
charged operator & 带电算符 \\
$p$-form (generalized) global symmetry & $p$-形式（广义）全局对称 \\
background gauge field & 背景规范场 \\
Wilson line; 't Hooft line & 威尔逊线；特霍夫特线 \\
electric / magnetic one-form symmetry & 电 / 磁 1-形式对称 \\
center symmetry & 中心对称 \\
spontaneous symmetry breaking (SSB) & 自发对称破缺 \\
Goldstone mode & 戈德斯通模 \\
Higgsing; confinement & Higgs 化；禁闭 \\
area law / perimeter law & 面积律 / 周长律 \\
Coleman--Mermin--Wagner theorem & CMW 定理 \\
't Hooft anomaly; anomaly matching & 特霍夫特反常；反常匹配 \\
BF theory & BF 理论 \\
higher-group (2-group, 3-group) symmetry & 高群对称 \\
non-invertible symmetry & 非可逆对称 \\
fusion rule & 融合规则 \\
selection rule & 选择定则 \\
\bottomrule
\end{tabular}
\end{center}

\vspace{1.2em}
\noindent\textit{声明：}本讲义为课前预习材料，不是 Brennan 教授的正式讲义。记号尽量靠近 arXiv:2306.00912，课堂若不同，以课堂为准。

\end{document}
